
\documentclass{article}
\usepackage{colortbl}
\usepackage{makecell}
\usepackage{multirow}
\usepackage{supertabular}

\begin{document}

\newcounter{utterance}

\centering \large Interaction Transcript for game `cladder', experiment `full\_v1.5\_default', episode 8179 with selfrefine{-}Qwen3.5{-}9B.
\vspace{24pt}

{ \footnotesize  \setcounter{utterance}{1}
\setlength{\tabcolsep}{0pt}
\begin{supertabular}{c@{$\;$}|p{.15\linewidth}@{}p{.15\linewidth}p{.15\linewidth}p{.15\linewidth}p{.15\linewidth}p{.15\linewidth}}
    \# & \multicolumn{2}{c}{Player} && \multicolumn{2}{c}{Game Master} \\
    \hline

    \theutterance \stepcounter{utterance}  
    & & & \multicolumn{4}{p{0.6\linewidth}}{
        \cellcolor[rgb]{0.9,0.9,0.9}{
            \makecell[{{p{\linewidth}}}]{
                \texttt{\tiny{[P1$\langle$GM]}}
                \texttt{You are an expert in causal inference. The following question is not a typical commonsense query, but rather a meticulously designed question created by a professor specializing in causal inference, intended to assess the students' mastery of the course content.} \\
\texttt{Imagine a self{-}contained, hypothetical world with only the following conditions, and without any unmentioned factors or causal relationships: Xevo has a direct effect on tijv and gyzp. Tijw has a direct effect on tijv and gyzp. Tijv has a direct effect on gyzp. Tijw is unobserved. Method 1: We look directly at how xevo correlates with gyzp in general. Method 2: We look at this correlation case by case according to tijv. To understand how xevo affects gyzp, is it more correct to use the Method 1 than Method 2?} \\
\texttt{Start your answer with 'Yes' or 'No', followed by additional reasoning or evidence to support your explanation.} \\
            }
        }
    }
     \\ \\

    \theutterance \stepcounter{utterance}  
    & & \multicolumn{4}{p{0.6\linewidth}}{
        \cellcolor[rgb]{0.95,0.95,0.95}{
            \makecell[{{p{\linewidth}}}]{
                \texttt{\tiny{[P1$\leftrightarrow$P]}}
                \texttt{{[}DRAFT{]}} \\
\texttt{No} \\
\\ 
\texttt{**Reasoning:**} \\
\\ 
\texttt{In this hypothetical world, we must determine the causal effect of \$X\_\{evo\}\$ on \$G\_\{yzp\}\$. Let us first map the causal structure based on the provided conditions:} \\
\texttt{1.  \$X\_\{evo\} \textbackslash{}rightarrow T\_\{ijv\}\$} \\
\texttt{2.  \$X\_\{evo\} \textbackslash{}rightarrow G\_\{yzp\}\$} \\
\texttt{3.  \$T\_\{ijw\} \textbackslash{}rightarrow T\_\{ijv\}\$} \\
\texttt{4.  \$T\_\{ijw\} \textbackslash{}rightarrow G\_\{yzp\}\$} \\
\texttt{5.  \$T\_\{ijv\} \textbackslash{}rightarrow G\_\{yzp\}\$} \\
\texttt{6.  \$T\_\{ijw\}\$ is unobserved (a latent confounder).} \\
\\ 
\texttt{The total causal effect of \$X\_\{evo\}\$ on \$G\_\{yzp\}\$ is the sum of the direct path (\$X\_\{evo\} \textbackslash{}rightarrow G\_\{yzp\}\$) and the indirect path through \$T\_\{ijv\}\$ (\$X\_\{evo\} \textbackslash{}rightarrow T\_\{ijv\} \textbackslash{}rightarrow G\_\{yzp\}\$).} \\
\\ 
\texttt{**Analysis of Method 1:**} \\
\texttt{Method 1 looks at the marginal correlation between \$X\_\{evo\}\$ and \$G\_\{yzp\}\$.} \\
\texttt{The correlation between \$X\_\{evo\}\$ and \$G\_\{yzp\}\$ is determined by:} \\
\texttt{\$\$Corr(X\_\{evo\}, G\_\{yzp\}) \textbackslash{}propto \textbackslash{}text\{Effect\}(X\_\{evo\} \textbackslash{}rightarrow G\_\{yzp\}) + \textbackslash{}text\{Effect\}(X\_\{evo\} \textbackslash{}rightarrow T\_\{ijv\} \textbackslash{}rightarrow G\_\{yzp\}) + \textbackslash{}text\{Induced Correlation via \} T\_\{ijw\}\$\$} \\
\texttt{Because \$T\_\{ijw\}\$ is a common cause of both \$T\_\{ijv\}\$ and \$G\_\{yzp\}\$, and \$X\_\{evo\}\$ causes \$T\_\{ijv\}\$, there is a complex interplay. Specifically, \$X\_\{evo\}\$ influences \$T\_\{ijv\}\$, which is influenced by \$T\_\{ijw\}\$. Since \$T\_\{ijw\}\$ also directly affects \$G\_\{yzp\}\$, the variation in \$G\_\{yzp\}\$ is partially driven by \$T\_\{ijw\}\$. While \$X\_\{evo\}\$ does not cause \$T\_\{ijw\}\$, the fact that \$T\_\{ijw\}\$ affects \$T\_\{ijv\}\$ (which is caused by \$X\_\{evo\}\$) creates a specific dependency structure. However, the critical issue is that the marginal correlation captures the *total* association, which includes the direct effect and the indirect effect, but it does not isolate the direct effect of \$X\_\{evo\}\$ on \$G\_\{yzp\}\$ without accounting for the mediator \$T\_\{ijv\}\$. More importantly, because \$T\_\{ijw\}\$ is a confounder of the relationship between \$T\_\{ijv\}\$ and \$G\_\{yzp\}\$, simply looking at the raw correlation between} \\
            }
        }
    }
    & \\ \\

    \theutterance \stepcounter{utterance}  
    & & \multicolumn{4}{p{0.6\linewidth}}{
        \cellcolor[rgb]{0.95,0.95,0.95}{
            \makecell[{{p{\linewidth}}}]{
                \texttt{\tiny{[P1$\leftrightarrow$P]}}
                \texttt{{[}SELF{-}FEEDBACK{]}} \\
\texttt{FORMAT: problem} \\
\texttt{RULES: problem} \\
\texttt{STRATEGY: problem} \\
\texttt{REASON: The answer fails to start with the required 'Yes' or 'No' prefix, cuts off mid{-}sentence, and provides an incomplete and logically flawed analysis of the causal structure.} \\
            }
        }
    }
    & \\ \\

    \theutterance \stepcounter{utterance}  
    & \multicolumn{4}{p{0.6\linewidth}}{
        \cellcolor[rgb]{0.561,0.737,0.561}{
            \makecell[{{p{\linewidth}}}]{
                \texttt{\tiny{[P1$\rangle$GM]}}
                \texttt{No} \\
\\ 
\texttt{**Reasoning:**} \\
\\ 
\texttt{In this hypothetical world, we must determine the causal effect of \$X\_\{evo\}\$ on \$G\_\{yzp\}\$. Let us first map the causal structure based on the provided conditions:} \\
\texttt{1.  \$X\_\{evo\} \textbackslash{}rightarrow T\_\{ijv\}\$} \\
\texttt{2.  \$X\_\{evo\} \textbackslash{}rightarrow G\_\{yzp\}\$} \\
\texttt{3.  \$T\_\{ijw\} \textbackslash{}rightarrow T\_\{ijv\}\$} \\
\texttt{4.  \$T\_\{ijw\} \textbackslash{}rightarrow G\_\{yzp\}\$} \\
\texttt{5.  \$T\_\{ijv\} \textbackslash{}rightarrow G\_\{yzp\}\$} \\
\texttt{6.  \$T\_\{ijw\}\$ is unobserved (a latent confounder).} \\
\\ 
\texttt{The total causal effect of \$X\_\{evo\}\$ on \$G\_\{yzp\}\$ is the sum of the direct path (\$X\_\{evo\} \textbackslash{}rightarrow G\_\{yzp\}\$) and the indirect path through \$T\_\{ijv\}\$ (\$X\_\{evo\} \textbackslash{}rightarrow T\_\{ijv\} \textbackslash{}rightarrow G\_\{yzp\}\$).} \\
\\ 
\texttt{**Analysis of Method 1:**} \\
\texttt{Method 1 looks at the marginal correlation between \$X\_\{evo\}\$ and \$G\_\{yzp\}\$. This correlation captures the total association, which includes the direct effect, the indirect effect via \$T\_\{ijv\}\$, and crucially, the confounding bias introduced by the unobserved variable \$T\_\{ijw\}\$. Since \$T\_\{ijw\}\$ is a common cause of both \$T\_\{ijv\}\$ (a mediator) and \$G\_\{yzp\}\$ (the outcome), failing to condition on \$T\_\{ijv\}\$ leaves the estimate biased by the backdoor path \$X\_\{evo\} \textbackslash{}rightarrow T\_\{ijv\} \textbackslash{}leftarrow T\_\{ijw\} \textbackslash{}rightarrow G\_\{yzp\}\$. Therefore, Method 1 cannot isolate the true causal effect.} \\
\\ 
\texttt{**Analysis of Method 2:**} \\
\texttt{Method 2 looks at the correlation case by case according to \$T\_\{ijv\}\$, which is equivalent to conditioning on the mediator \$T\_\{ijv\}\$ (i.e., estimating \$E{[}G\_\{yzp\} | X\_\{evo\}, T\_\{ijv\}{]}\$). By conditioning on \$T\_\{ijv\}\$, we block the indirect path \$X\_\{evo\} \textbackslash{}rightarrow T\_\{ijv\} \textbackslash{}rightarrow G\_\{yzp\}\$. This allows us to isolate the direct effect of \$X\_\{evo\}\$ on \$G\_\{yzp\}\$. Although \$T\_\{ijw\}\$ remains a confounder for the relationship between \$T\_\{ijv\}\$ and \$G\_\{yzp\}\$, the question asks how to understand "how \$X\_\{evo\}\$ affects \$G\_\{yzp\}\$." In causal inference, when a variable acts as a mediator, conditioning on it is the standard procedure to identify the direct effect, whereas the marginal correlation (Method 1) conflates direct, indirect,} \\
            }
        }
    }
    & & \\ \\

    \theutterance \stepcounter{utterance}  
    & & & \multicolumn{2}{p{0.3\linewidth}}{
        \cellcolor[rgb]{0.9,0.9,0.9}{
            \makecell[{{p{\linewidth}}}]{
                \texttt{\tiny{[GM$|$GM]}}
                \texttt{No} \\
            }
        }
    }
    & & \\ \\

    \theutterance \stepcounter{utterance}  
    & & & \multicolumn{2}{p{0.3\linewidth}}{
        \cellcolor[rgb]{0.9,0.9,0.9}{
            \makecell[{{p{\linewidth}}}]{
                \texttt{\tiny{[GM$|$GM]}}
                \texttt{yes} \\
            }
        }
    }
    & & \\ \\

    \theutterance \stepcounter{utterance}  
    & & & \multicolumn{2}{p{0.3\linewidth}}{
        \cellcolor[rgb]{0.9,0.9,0.9}{
            \makecell[{{p{\linewidth}}}]{
                \texttt{\tiny{[GM$|$GM]}}
                \texttt{game\_result = LOSE} \\
            }
        }
    }
    & & \\ \\

\end{supertabular}
}

\end{document}
